\section*{Заключение}
\addcontentsline{toc}{section}{Заключение}

В ходе выполнения курсовой работы была рассмотрена предметная область интернет-торговли
цифровой электроникой и определены требования к веб-приложению для управления ассортиментом
и заказами. Анализ существующих решений показал, что крупные интернет-магазины и
маркетплейсы обладают развитым функционалом, однако не всегда подходят для самостоятельного
магазина, которому требуется собственная система управления товарами, клиентами и заказами.

По результатам работы можно сформулировать следующие \textbf{выводы}:

\begin{enumerate}
    \item Для магазина цифровой электроники ключевыми функциями являются каталог с
    характеристиками товаров, поиск, фильтрация, корзина, оформление заказа и история покупок.
    \item Для серверной части проекта целесообразно использовать Django, так как он предоставляет
    встроенные средства работы с базой данных, административной панелью, пользователями и правами
    доступа.
    \item Модель данных должна учитывать связи между пользователями, категориями, товарами,
    корзиной, заказами и позициями заказа, чтобы обеспечить корректную обработку покупок.
    \item Разграничение ролей покупателей, менеджеров и администраторов повышает безопасность
    системы и упрощает организацию рабочих процессов магазина.
\end{enumerate}

\textbf{Перспективы развития проекта} связаны с добавлением онлайн-оплаты, интеграции со
службами доставки, системы отзывов, промокодов, уведомлений о статусе заказа и аналитики продаж.

При разработке пояснительной записки к курсовому проекту соблюдались требования~\cite{gost_report},
определяющие структуру и правила оформления научно-исследовательских работ.
